% Options for packages loaded elsewhere
\PassOptionsToPackage{unicode}{hyperref}
\PassOptionsToPackage{hyphens}{url}
\documentclass[
]{article}
\usepackage{xcolor}
\usepackage[margin=1in]{geometry}
\usepackage{amsmath,amssymb}
\setcounter{secnumdepth}{-\maxdimen} % remove section numbering
\usepackage{iftex}
\ifPDFTeX
  \usepackage[T1]{fontenc}
  \usepackage[utf8]{inputenc}
  \usepackage{textcomp} % provide euro and other symbols
\else % if luatex or xetex
  \usepackage{unicode-math} % this also loads fontspec
  \defaultfontfeatures{Scale=MatchLowercase}
  \defaultfontfeatures[\rmfamily]{Ligatures=TeX,Scale=1}
\fi
\usepackage{lmodern}
\ifPDFTeX\else
  % xetex/luatex font selection
\fi
% Use upquote if available, for straight quotes in verbatim environments
\IfFileExists{upquote.sty}{\usepackage{upquote}}{}
\IfFileExists{microtype.sty}{% use microtype if available
  \usepackage[]{microtype}
  \UseMicrotypeSet[protrusion]{basicmath} % disable protrusion for tt fonts
}{}
\makeatletter
\@ifundefined{KOMAClassName}{% if non-KOMA class
  \IfFileExists{parskip.sty}{%
    \usepackage{parskip}
  }{% else
    \setlength{\parindent}{0pt}
    \setlength{\parskip}{6pt plus 2pt minus 1pt}}
}{% if KOMA class
  \KOMAoptions{parskip=half}}
\makeatother
\usepackage{color}
\usepackage{fancyvrb}
\newcommand{\VerbBar}{|}
\newcommand{\VERB}{\Verb[commandchars=\\\{\}]}
\DefineVerbatimEnvironment{Highlighting}{Verbatim}{commandchars=\\\{\}}
% Add ',fontsize=\small' for more characters per line
\usepackage{framed}
\definecolor{shadecolor}{RGB}{248,248,248}
\newenvironment{Shaded}{\begin{snugshade}}{\end{snugshade}}
\newcommand{\AlertTok}[1]{\textcolor[rgb]{0.94,0.16,0.16}{#1}}
\newcommand{\AnnotationTok}[1]{\textcolor[rgb]{0.56,0.35,0.01}{\textbf{\textit{#1}}}}
\newcommand{\AttributeTok}[1]{\textcolor[rgb]{0.13,0.29,0.53}{#1}}
\newcommand{\BaseNTok}[1]{\textcolor[rgb]{0.00,0.00,0.81}{#1}}
\newcommand{\BuiltInTok}[1]{#1}
\newcommand{\CharTok}[1]{\textcolor[rgb]{0.31,0.60,0.02}{#1}}
\newcommand{\CommentTok}[1]{\textcolor[rgb]{0.56,0.35,0.01}{\textit{#1}}}
\newcommand{\CommentVarTok}[1]{\textcolor[rgb]{0.56,0.35,0.01}{\textbf{\textit{#1}}}}
\newcommand{\ConstantTok}[1]{\textcolor[rgb]{0.56,0.35,0.01}{#1}}
\newcommand{\ControlFlowTok}[1]{\textcolor[rgb]{0.13,0.29,0.53}{\textbf{#1}}}
\newcommand{\DataTypeTok}[1]{\textcolor[rgb]{0.13,0.29,0.53}{#1}}
\newcommand{\DecValTok}[1]{\textcolor[rgb]{0.00,0.00,0.81}{#1}}
\newcommand{\DocumentationTok}[1]{\textcolor[rgb]{0.56,0.35,0.01}{\textbf{\textit{#1}}}}
\newcommand{\ErrorTok}[1]{\textcolor[rgb]{0.64,0.00,0.00}{\textbf{#1}}}
\newcommand{\ExtensionTok}[1]{#1}
\newcommand{\FloatTok}[1]{\textcolor[rgb]{0.00,0.00,0.81}{#1}}
\newcommand{\FunctionTok}[1]{\textcolor[rgb]{0.13,0.29,0.53}{\textbf{#1}}}
\newcommand{\ImportTok}[1]{#1}
\newcommand{\InformationTok}[1]{\textcolor[rgb]{0.56,0.35,0.01}{\textbf{\textit{#1}}}}
\newcommand{\KeywordTok}[1]{\textcolor[rgb]{0.13,0.29,0.53}{\textbf{#1}}}
\newcommand{\NormalTok}[1]{#1}
\newcommand{\OperatorTok}[1]{\textcolor[rgb]{0.81,0.36,0.00}{\textbf{#1}}}
\newcommand{\OtherTok}[1]{\textcolor[rgb]{0.56,0.35,0.01}{#1}}
\newcommand{\PreprocessorTok}[1]{\textcolor[rgb]{0.56,0.35,0.01}{\textit{#1}}}
\newcommand{\RegionMarkerTok}[1]{#1}
\newcommand{\SpecialCharTok}[1]{\textcolor[rgb]{0.81,0.36,0.00}{\textbf{#1}}}
\newcommand{\SpecialStringTok}[1]{\textcolor[rgb]{0.31,0.60,0.02}{#1}}
\newcommand{\StringTok}[1]{\textcolor[rgb]{0.31,0.60,0.02}{#1}}
\newcommand{\VariableTok}[1]{\textcolor[rgb]{0.00,0.00,0.00}{#1}}
\newcommand{\VerbatimStringTok}[1]{\textcolor[rgb]{0.31,0.60,0.02}{#1}}
\newcommand{\WarningTok}[1]{\textcolor[rgb]{0.56,0.35,0.01}{\textbf{\textit{#1}}}}
\usepackage{graphicx}
\makeatletter
\newsavebox\pandoc@box
\newcommand*\pandocbounded[1]{% scales image to fit in text height/width
  \sbox\pandoc@box{#1}%
  \Gscale@div\@tempa{\textheight}{\dimexpr\ht\pandoc@box+\dp\pandoc@box\relax}%
  \Gscale@div\@tempb{\linewidth}{\wd\pandoc@box}%
  \ifdim\@tempb\p@<\@tempa\p@\let\@tempa\@tempb\fi% select the smaller of both
  \ifdim\@tempa\p@<\p@\scalebox{\@tempa}{\usebox\pandoc@box}%
  \else\usebox{\pandoc@box}%
  \fi%
}
% Set default figure placement to htbp
\def\fps@figure{htbp}
\makeatother
\setlength{\emergencystretch}{3em} % prevent overfull lines
\providecommand{\tightlist}{%
  \setlength{\itemsep}{0pt}\setlength{\parskip}{0pt}}
\usepackage{bookmark}
\IfFileExists{xurl.sty}{\usepackage{xurl}}{} % add URL line breaks if available
\urlstyle{same}
\hypersetup{
  pdftitle={Discrete distributions: Bernoulli{,} binomial{,} Poisson},
  hidelinks,
  pdfcreator={LaTeX via pandoc}}

\title{Discrete distributions: Bernoulli, binomial, Poisson}
\author{}
\date{\vspace{-2.5em}}

\begin{document}
\maketitle

\textbf{Inference labs} use the five-step template: Hypothesis →
Visualise → Assumptions → Conduct → Conclude.

\subsection{Learning objectives}\label{learning-objectives}

\begin{itemize}
\tightlist
\item
  Simulate Bernoulli trials and aggregate them to a binomial.
\item
  Show the Poisson approximation to the binomial when \emph{n} is large
  and \emph{p} is small.
\item
  Use \texttt{dbinom()}, \texttt{pbinom()}, and \texttt{rbinom()} (and
  their Poisson counterparts) to answer probability questions.
\end{itemize}

\subsection{Prerequisites}\label{prerequisites}

Lab 2.2 (probability).

\subsection{Background}\label{background}

The Bernoulli distribution is the simplest non-trivial distribution: one
coin, two outcomes, a single probability \emph{p}. The binomial adds
structure --- \emph{n} independent Bernoulli trials with the same
\emph{p} --- and gives the distribution of the count of successes. The
Poisson arises in the limit where \emph{n} is large, \emph{p} is small,
and \emph{np} stays fixed; it describes counts of rare events over a
fixed window of opportunity, without needing to specify \emph{n} and
\emph{p} separately.

Most count-like quantities in biomedicine can be modelled by one of
these three. Whether a tumour is responding (Bernoulli), how many of
twenty patients relapse (binomial), how many emergency admissions arrive
in an hour (Poisson) --- the same three distributions cover a great deal
of ground. The simulation-first approach makes them concrete before we
write a formula.

There is a common conflation between the binomial and the Poisson in the
clinical literature. When the sample size is large and the event is
rare, the two distributions give nearly identical probabilities and
either is defensible. When the event is not rare, the Poisson can assign
meaningful probability to impossible events (more successes than
trials), so the binomial is strictly better.

\subsection{Setup}\label{setup}

\begin{Shaded}
\begin{Highlighting}[]
\FunctionTok{library}\NormalTok{(tidyverse)}
\FunctionTok{set.seed}\NormalTok{(}\DecValTok{42}\NormalTok{)}
\FunctionTok{theme\_set}\NormalTok{(}\FunctionTok{theme\_minimal}\NormalTok{(}\AttributeTok{base\_size =} \DecValTok{12}\NormalTok{))}
\end{Highlighting}
\end{Shaded}

\subsection{1. Hypothesis}\label{hypothesis}

We characterise three discrete distributions by simulation, then compare
the empirical frequencies to the theoretical probabilities. No
hypothesis is tested.

\subsection{2. Visualise}\label{visualise}

A Bernoulli trial: one coin, success probability 0.3.

\begin{Shaded}
\begin{Highlighting}[]
\NormalTok{bern }\OtherTok{\textless{}{-}} \FunctionTok{rbinom}\NormalTok{(n\_trials, }\DecValTok{1}\NormalTok{, }\FloatTok{0.3}\NormalTok{)}
\FunctionTok{mean}\NormalTok{(bern)   }\CommentTok{\# empirical p}
\end{Highlighting}
\end{Shaded}

A binomial count: how many successes in 20 trials each at \emph{p} =
0.3, repeated 10,000 times.

\begin{Shaded}
\begin{Highlighting}[]
\NormalTok{bin  }\OtherTok{\textless{}{-}} \FunctionTok{rbinom}\NormalTok{(n\_rep, }\AttributeTok{size =} \DecValTok{20}\NormalTok{, }\AttributeTok{prob =} \FloatTok{0.3}\NormalTok{)}

\NormalTok{bin\_tibble }\OtherTok{\textless{}{-}} \FunctionTok{tibble}\NormalTok{(}\AttributeTok{k =}\NormalTok{ bin) }\SpecialCharTok{|\textgreater{}}
  \FunctionTok{count}\NormalTok{(k) }\SpecialCharTok{|\textgreater{}}
  \FunctionTok{mutate}\NormalTok{(}\AttributeTok{prop =}\NormalTok{ n }\SpecialCharTok{/} \FunctionTok{sum}\NormalTok{(n),}
         \AttributeTok{theo =} \FunctionTok{dbinom}\NormalTok{(k, }\DecValTok{20}\NormalTok{, }\FloatTok{0.3}\NormalTok{))}
\end{Highlighting}
\end{Shaded}

\begin{Shaded}
\begin{Highlighting}[]
\NormalTok{bin\_tibble }\SpecialCharTok{|\textgreater{}}
  \FunctionTok{ggplot}\NormalTok{(}\FunctionTok{aes}\NormalTok{(k, prop)) }\SpecialCharTok{+}
  \FunctionTok{geom\_col}\NormalTok{(}\AttributeTok{fill =} \StringTok{"grey70"}\NormalTok{, }\AttributeTok{alpha =} \FloatTok{0.8}\NormalTok{) }\SpecialCharTok{+}
  \FunctionTok{geom\_point}\NormalTok{(}\FunctionTok{aes}\NormalTok{(}\AttributeTok{y =}\NormalTok{ theo), }\AttributeTok{colour =} \StringTok{"firebrick"}\NormalTok{, }\AttributeTok{size =} \DecValTok{2}\NormalTok{) }\SpecialCharTok{+}
  \FunctionTok{labs}\NormalTok{(}\AttributeTok{x =} \StringTok{"Successes in 20 trials"}\NormalTok{,}
       \AttributeTok{y =} \StringTok{"Proportion (bars) vs. P(k) (points)"}\NormalTok{)}
\end{Highlighting}
\end{Shaded}

\subsection{3. Assumptions}\label{assumptions}

Binomial assumes fixed \emph{n}, independence, and constant \emph{p}.
Poisson assumes a constant rate and independent counts over disjoint
intervals. When \emph{n} is large and \emph{p} is small with \emph{np} =
\emph{λ}, the binomial approaches Poisson.

\begin{Shaded}
\begin{Highlighting}[]
\NormalTok{probs }\OtherTok{\textless{}{-}} \FunctionTok{tibble}\NormalTok{(}
  \AttributeTok{k =} \DecValTok{0}\SpecialCharTok{:}\DecValTok{10}\NormalTok{,}
  \AttributeTok{binom  =} \FunctionTok{dbinom}\NormalTok{(k, }\DecValTok{1000}\NormalTok{, }\FloatTok{0.003}\NormalTok{),}
  \AttributeTok{pois   =} \FunctionTok{dpois}\NormalTok{(k, }\DecValTok{3}\NormalTok{)}
\NormalTok{)}
\NormalTok{probs}
\end{Highlighting}
\end{Shaded}

\subsection{4. Conduct}\label{conduct}

Answer two specific questions.

\begin{Shaded}
\begin{Highlighting}[]
\CommentTok{\# P(at least 8 responders)?}
\NormalTok{p\_ge8 }\OtherTok{\textless{}{-}} \DecValTok{1} \SpecialCharTok{{-}} \FunctionTok{pbinom}\NormalTok{(}\DecValTok{7}\NormalTok{, }\AttributeTok{size =} \DecValTok{20}\NormalTok{, }\AttributeTok{prob =} \FloatTok{0.3}\NormalTok{)}
\NormalTok{p\_ge8}

\CommentTok{\# Q2: An ER admits patients at a rate of 3/hour. What is}
\CommentTok{\# P(7 or more patients in a given hour)?}
\NormalTok{p\_pois\_ge7 }\OtherTok{\textless{}{-}} \DecValTok{1} \SpecialCharTok{{-}} \FunctionTok{ppois}\NormalTok{(}\DecValTok{6}\NormalTok{, }\AttributeTok{lambda =} \DecValTok{3}\NormalTok{)}
\NormalTok{p\_pois\_ge7}
\end{Highlighting}
\end{Shaded}

Simulate Q2 to confirm.

\begin{Shaded}
\begin{Highlighting}[]
\FunctionTok{mean}\NormalTok{(sim }\SpecialCharTok{\textgreater{}=} \DecValTok{7}\NormalTok{)}
\end{Highlighting}
\end{Shaded}

\begin{Shaded}
\begin{Highlighting}[]
\FunctionTok{tibble}\NormalTok{(}\AttributeTok{k =} \DecValTok{0}\SpecialCharTok{:}\DecValTok{15}\NormalTok{) }\SpecialCharTok{|\textgreater{}}
  \FunctionTok{mutate}\NormalTok{(}\AttributeTok{pois\_density =} \FunctionTok{dpois}\NormalTok{(k, }\DecValTok{3}\NormalTok{)) }\SpecialCharTok{|\textgreater{}}
  \FunctionTok{ggplot}\NormalTok{(}\FunctionTok{aes}\NormalTok{(k, pois\_density)) }\SpecialCharTok{+}
  \FunctionTok{geom\_col}\NormalTok{(}\AttributeTok{fill =} \StringTok{"steelblue"}\NormalTok{) }\SpecialCharTok{+}
  \FunctionTok{labs}\NormalTok{(}\AttributeTok{x =} \StringTok{"Number of events"}\NormalTok{, }\AttributeTok{y =} \StringTok{"P(X = k)"}\NormalTok{)}
\end{Highlighting}
\end{Shaded}

\subsection{5. Concluding statement}\label{concluding-statement}

\begin{quote}
A response rate of 0.3 in a 20-patient trial implies a probability of
\texttt{round(p\_ge8,\ 3)} of observing eight or more responders. An ER
with a mean admission rate of 3/hour has a probability of
\texttt{round(p\_pois\_ge7,\ 3)} of receiving seven or more admissions
in a given hour; 10,000 simulated hours recovered an empirical
proportion of \texttt{round(mean(sim\ \textgreater{}=\ 7),\ 3)}. The
binomial with \emph{n} = 1000, \emph{p} = 0.003 and the Poisson with λ =
3 gave nearly indistinguishable probabilities across \emph{k} = 0--10.
\end{quote}

Discrete distributions are under-appreciated. Many designs that are
reached for as continuous (eg regressions on a rate) have a cleaner
count-based model sitting next to them.

\subsection{Common pitfalls}\label{common-pitfalls}

\begin{itemize}
\tightlist
\item
  Using the normal approximation to the binomial when \emph{n} is small.
  For small \emph{n} use \texttt{pbinom()} directly.
\item
  Modelling a rate as a proportion when the denominator is person-time,
  not persons.
\item
  Treating a Poisson with a known \emph{λ} as if its variance had to be
  estimated. For Poisson, mean = variance.
\end{itemize}

\subsection{Further reading}\label{further-reading}

\begin{itemize}
\tightlist
\item
  Rosner B. \emph{Fundamentals of Biostatistics}, chapter on discrete
  distributions.
\item
  Agresti A. \emph{Categorical Data Analysis}, opening chapter.
\end{itemize}

\subsection{Session info}\label{session-info}

\begin{Shaded}
\begin{Highlighting}[]

\end{Highlighting}
\end{Shaded}

\subsection{See also --- chapter index}\label{see-also-chapter-index}

\begin{itemize}
\tightlist
\item
  \href{../index.Rmd}{Methods in this chapter}
\item
  \href{../../../ch00-find-your-method.Rmd}{Find your method (Chapter
  0)}
\end{itemize}

\end{document}
